\documentclass[journal=jacsat,manuscript=article]{achemso}
\usepackage{amsmath}
\usepackage{amssymb}
\usepackage{graphicx}
\usepackage{orcidlink}
\usepackage{adjustbox}
\usepackage{multicol}
\usepackage{multirow}
\usepackage{mathtools}
\usepackage[version=4]{mhchem}
\usepackage{booktabs}
\usepackage{xcolor}
\usepackage{pgfplots}
\pgfplotsset{compat=1.18}
\usepackage{tikz}
\usepackage{fontawesome5}
\usetikzlibrary{positioning,arrows.meta,shapes.geometric}
\usepackage{bm}

\newcommand{\vect}[1]{\bm{#1}}
\newcommand{\tensor}[1]{\bm{#1}}
\newcommand{\rank}{\operatorname{rank}}

\title{Representation Equivalence of Centrifugal Distortion Constants: Linear Transport and the First Breakdown}

\author{Vincenzo Barone\,\orcidlink{0000-0001-6420-4107}}
\affiliation[c]{INSTM, via G. Giusti 9, 50121 Firenze, Italy}
\email{vincebarone52@gmail.com}

\begin{document}

\begin{tocentry}
%\includegraphics[width=\textwidth]{Figures/toc_ebn.png}
\end{tocentry}

\begin{abstract}
Quartic and sextic centrifugal distortion constants are usually handled as
representation-dependent parameters of Watson's effective rotational
Hamiltonian, whether they are obtained from spectroscopic fits or from direct
quantum-chemical calculations. In this work their representation equivalence
is reformulated in terms of linear tensor transport and its first breakdown.
The standard quartic and standard sextic sectors are recast in a common
tensorial framework in which Watson constants are treated as projected
coordinates of underlying tensor representatives reconstructed by
Moore--Penrose pseudoinversion, which thus acts as a natural gauge fixing of
the inverse problem.

Within this formulation, standard quartic and standard sextic centrifugal
distortion are shown to belong to the same linear
representation-equivalent sector. The key structural point is that the standard
sextic geometry sector can be written in the same first-level tensorial
language that governs the standard quartic problem, so that second derivatives
of the inertia tensor need not be treated as primitive objects. The explicit
anharmonic sextic terms then enter only through the cubic force constants.
Representation and reduction changes are therefore performed by lifting Watson
constants to tensor representatives, transporting the tensors under axis
permutations, and reprojecting them onto the target Hamiltonian
parametrization.

Quartic and standard sextic centrifugal distortion constants are thus shown to
belong to a common linear tensor sector that can be transported between
representations and validated at essentially harmonic computational cost. When
genuine three-index cubic contributions are small, the full set of distortion
constants can be reconstructed at roughly twice the cost of a harmonic
calculation using analytic gradients, with geometric, harmonic, and cubic
contributions evaluated independently at different levels of theory.

This common standard sector provides the natural baseline from which the first
departure from linear representation equivalence can be identified. From the
quantum-chemical point of view, the most important beyond-standard term is in
general the numerically largest one. From the spectroscopic fitting point of
view, however, the crucial term is the first one that breaks linear transport
between representations. On the quartic side, the channel $H_{22}$ is shown to
be the first contribution of this type, because it requires a genuinely new
intrinsic tensorial object beyond the standard first-level structure. On the
sextic side, the corresponding first post-standard candidate is the term
linear in $\tau(H_{22})$ inside the sextic geometry functional. These
quantities are not introduced as complete higher-order corrections, but as
low-cost tensorial diagnostics of both the sensitivity of fitted constants to
representation/reduction choices and the need to move beyond the standard
perturbative sector. The main text is therefore focused on linear
representation transport and its first breakdown, while the detailed tensor
constructions and derivations that justify this structure are collected in the
appendices.
\end{abstract}

\section{Introduction}

High--resolution rotational spectroscopy provides extremely precise
information on molecular structure and rovibrational dynamics.
\cite{TownesSchawlow1955,GordyCook1984,BunkerJensen2005}
Both experiment and theory ultimately access the same physical
observables, namely rotational transition frequencies. Their
interpretation, however, naturally leads to two complementary
viewpoints. In the \emph{inverse} problem, spectroscopic parameters are
determined from measured spectra by fitting an effective Hamiltonian
to observed transition frequencies.
\cite{GordyCook1984,TownesSchawlow1955}
In the corresponding \emph{direct} problem, the same parameters are
predicted from a molecular potential energy surface and the associated
rovibrational couplings.\cite{Dinu2022}

In practical spectroscopy, rotational spectra are commonly analyzed
using Watson's effective rotational Hamiltonian,
\cite{Watson1968,Watson1977}
which includes quartic and sextic centrifugal distortion terms.
Within this phenomenological framework one works with five quartic and
seven sextic centrifugal distortion constants, which are typically
treated as adjustable parameters in least--squares fits and
implemented in widely used programs such as SPFIT/SPCAT and
PGOPHER.\cite{Pickett1991,Western2017}
At the same time, modern quantum-chemical and rovibrational methods
can access the same constants from direct calculations. The central
question is therefore no longer whether centrifugal distortion
constants belong to experiment or to theory, but how the same
underlying rovibrational information should be represented and
compared across different Hamiltonian reductions, principal-axis
representations, and computational schemes.

A basic complication is that Watson's centrifugal distortion constants
are not intrinsic observables. Their numerical values depend on the
chosen reduction and on the adopted principal-axis representation.
Different parameter sets may therefore describe the same underlying
rovibrational content equally well. In this sense, the spectroscopic
constants are best regarded as coordinates of an effective Hamiltonian
parametrization rather than as primary molecular descriptors.

This viewpoint suggests that the natural comparison level is not the
Watson constants themselves, but suitable tensor representatives from
which those constants are obtained by projection. In the quartic case,
the standard reduced object is the symmetric tensor $\tau'$. In the
sextic case, the corresponding physical content is represented by a
canonical reduced sector rather than by a single tensor of quartic
type. Once this tensorial level is adopted, representation and
reduction changes may be formulated as transport operations on the
tensor representatives, followed by reprojection onto the desired
spectroscopic parametrization.

An important structural point then emerges immediately. In the quartic
case, the tensor representation has one more degree of freedom than
the five spectroscopic constants of Watson's Hamiltonian. The inverse
reconstruction of the tensor from the constants is therefore not
one-to-one, but defines an affine gauge problem. In the present work
this inverse step is fixed by the Moore--Penrose pseudoinverse, which
provides a natural and numerically stable gauge choice. The
pseudoinverse does not impose a physical constraint on centrifugal
distortion; it selects one representative within the family of
tensorial objects compatible with a given spectroscopic parameter set.

The same logic extends naturally to the sextic problem. However, the
main point of the present work is not to build the most general tensor
theory of quartic and sextic centrifugal distortion. Rather, the goal
is to identify the \emph{standard linear sector} in which quartic and
sextic centrifugal distortion constants admit a common tensorial
description and transform linearly between representations after
pseudoinverse gauge fixing. This is the sector that underlies standard
Watson-style analyses and the usual VPT2 treatment of centrifugal
distortion.

In the present work the perturbative bookkeeping follows the standard
Watson/VPT2-style partition. The labels $H_{nm}$ indicate the
vibrational and rotational bidegree of each block, but do not by
themselves fix its perturbative order. Within the partition adopted
here, $H_{03}$, $H_{12}$, and $H_{30}$ are first-order terms, whereas
$H_{40}$, $H_{22}$, $H_{13}$, $H_{31}$, and $H_{04}$ are second-order
terms. This convention is standard in the present context, but it is
not universal: in alternative formulations, including state-summation
schemes that are also often labeled as VPT2, the effective
perturbative order assigned to the same blocks may differ.

Within this standard partition, the quartic and sextic constants are
not treated here as unrelated sectors. A central result of the present
analysis is that the standard sextic structure closes on the same
first-level tensorial language that governs the standard quartic
problem. In particular, the standard sextic geometry sector can be
written without taking second derivatives of the inertia tensor as
primitive objects. The explicit anharmonic sextic contribution then
enters only through the cubic force constants. This identifies a
common standard linear sector for quartic and sextic centrifugal
distortion.

This standard sector is the natural place where direct and inverse
descriptions become strictly comparable. Spectroscopic constants may
be lifted by pseudoinverse reconstruction to tensor representatives,
transported linearly under axis permutations, and reprojected in any
desired reduction or representation. In this sense, standard quartic
and standard sextic centrifugal distortion constants belong to one
common representation-equivalent framework.

The second goal of the present work is to identify the first departure
from this standard linear sector. On the quartic side, this role is
played by the channel $H_{22}$, which is the first contribution that
requires a genuinely new intrinsic tensorial object beyond the
standard first-level structure. On the sextic side, the corresponding
first post-standard candidate is the term linear in $\tau(H_{22})$
inside the sextic geometry functional. These quantities are not
introduced as complete higher-order corrections. Rather, they are used
as low-cost diagnostics of when the standard linear sector begins to
break down, and therefore of when one should expect a stronger
dependence on representation choice or the need to move beyond the
standard perturbative level.

From this perspective, the purpose of the present paper is twofold.
First, it establishes a unified tensorial framework for the transport
of standard quartic and standard sextic centrifugal distortion
constants between different reductions and principal-axis
representations. Second, it identifies the first quartic and sextic
diagnostics that signal the breakdown of this linear standard sector.
The corresponding transformations and diagnostics are made available
through practical tools, either as a macOS app or as a Python code,
so that the framework can be used directly on fitted or computed
parameter sets.
The detailed tensor constructions and derivations that justify this
structure are collected in the appendices, while the main text is
focused on the transport problem itself and on its first breakdown.

The paper is organized as follows. Section~2 introduces the adopted
perturbative framework, the tensor representatives, and the
pseudoinverse gauge-fixing procedure. Section~3 discusses the
standard quartic sector. Section~4 treats the standard sextic sector
and shows how it closes on the same first-level tensorial language.
Section~5 identifies the common standard linear sector. Section~6
introduces $H_{22}$ as the first quartic breakdown term, and
Section~7 discusses the corresponding first sextic analogue, namely
the term linear in $\tau(H_{22})$. The appendices collect the
detailed tensor derivations, projection formulas, and technical
implementation details.

\bibliography{CeDiTT3}

\end{document}
